(1) 由给定色散关系取 \(\omega(k)=\sqrt{gk+ak^3}\)，\(a=\sigma/\rho\)。相速度 \(c_p=\omega/k=\sqrt{g/k+ak}\)；群速度 \(c_g=d\omega/dk=(g+3ak^2)/[2\sqrt{gk+ak^3}]\)。两者以不同的 \(k\) 函数变化，故波包与单个波峰不会始终同步。
(2) 最小化相速度的平方 \(c_p^2=g/k+ak\)，求导 \(-g/k^2+a=0\)，得 \(k_*=\sqrt{g/a}=\sqrt{\rho g/\sigma}\)。二阶导 \(2g/k^3>0\)，确为唯一最小值。于是 \(c_{p,\min}^2=g/k_*+ak_*=2\sqrt{ag}\)，故 \(c_{p,\min}=(4\sigma g/\rho)^{1/4}\)。
(3) 当 \(k\ll k_*\) 时重力项占优，\(\omega\sim\sqrt{gk}\)，于是 \(c_g/c_p\to1/2\)，长波波包比相位传播慢；当 \(k\gg k_*\) 时毛细项占优，\(\omega\sim\sqrt a\,k^{3/2}\)，所以 \(c_g/c_p\to3/2\)，高波数波包反而更快。
(4) 在极小相速度处因 \(c_p'(k_*)=0\)，有 \(c_g=d(kc_p)/dk=c_p+kc_p'=c_p\)。深水条件为 \(kh\gg1\)，小振幅要求波陡度 \(k\eta_0\ll1\)，还需忽略粘性耗散；浅水、强非线性或显著粘性都会改变既定色散曲线和速度排序。
